\documentclass[11pt]{article}
\usepackage[utf8]{inputenc}
\usepackage[T1]{fontenc}
\usepackage[spanish]{babel}
\usepackage{amsmath}
\usepackage{hyperref}

\title{Artículo simple}
\author{Fixture de LaTeX Viewer}
\date{}

\begin{document}
\maketitle

\section{Introducción}\label{sec:intro}

Este es un párrafo largo escrito en una sola línea, como en el libro de prueba: no hay cortes manuales y el párrafo entero vive en la misma línea del archivo fuente, de modo que saltar «a la línea» desde el PDF no alcanza y hace falta caer en la palabra exacta. Sirve para probar el refinamiento por palabra de SyncTeX, la búsqueda inversa y el soft wrap del editor con una línea de más de trescientos caracteres que se reparte en varias líneas del PDF.

Un segundo párrafo, también de una sola línea, que cita la ecuación~\eqref{eq:cuadratica} y la sección~\ref{sec:resultado}. La primera referencia obliga a una segunda pasada de LaTeX, que es lo que se mide en las líneas de base.

\section{Resultado}\label{sec:resultado}

La ecuación cuadrática es
\begin{equation}\label{eq:cuadratica}
  x = \frac{-b \pm \sqrt{b^2 - 4ac}}{2a}.
\end{equation}

Cierra con otro párrafo largo para que haya texto después de la ecuación: se vuelve a la sección~\ref{sec:intro} y se comprueba que el número de página y el de sección salen bien resueltos sin signos de interrogación.

\end{document}
